% META: {"id":"1SPE-PRODSCAL-CO-009","chapitre":"1SPE-PRODUIT-SCALAIRE","type_objet":"corrige","exercice_id":"1SPE-PRODSCAL-EX-009","status":"generated"}
% BEGIN-VERIFY
% from sympy import Matrix, sqrt, Rational, simplify, acos, pi, N
% A, B, C = Matrix([0,0]), Matrix([6,0]), Matrix([2,4])
% AB, AC, BC = B-A, C-A, C-B
% CA, CB = A-C, B-C
% assert AB.dot(AC) == 12
% assert AB.dot(BC) == -24
% assert AC.dot(BC) == 8
% assert CA.dot(CB) == 8
% cos_A = AB.dot(AC)/(AB.norm()*AC.norm())
% cos_B = (-AB).dot(BC)/(AB.norm()*BC.norm())
% cos_C = CA.dot(CB)/(CA.norm()*CB.norm())
% assert simplify(cos_A - 1/sqrt(5)) == 0
% assert simplify(cos_B - 1/sqrt(2)) == 0
% assert simplify(cos_C - 1/sqrt(10)) == 0
% assert all(0 < value < 1 for value in (cos_A,cos_B,cos_C))
% angle_A, angle_B, angle_C = [acos(value) for value in (cos_A,cos_B,cos_C)]
% assert angle_B == pi/4
% # The angle sum is exact by the triangle-angle theorem; this checks the numerical calculation.
% assert abs(N(angle_A+angle_B+angle_C-pi,30)) < Rational(1,10**25)
% assert Rational(6335,100) < N(angle_A*180/pi,20) < Rational(6345,100)
% assert Rational(7155,100) < N(angle_C*180/pi,20) < Rational(7165,100)
% END-VERIFY
\begin{corrige}{1SPE-PRODSCAL-EX-009}
\begin{enumerate}
\item $\overrightarrow{AB}\begin{pmatrix} 6 \\ 0 \end{pmatrix}$, $\overrightarrow{AC}\begin{pmatrix} 2 \\ 4 \end{pmatrix}$, $\overrightarrow{BC}\begin{pmatrix} -4 \\ 4 \end{pmatrix}$.

$\overrightarrow{AB} \cdot \overrightarrow{AC} = 12$, $\overrightarrow{AB} \cdot \overrightarrow{BC} = -24$, $\overrightarrow{AC} \cdot \overrightarrow{BC} = -8 + 16 = 8$.

\item $\|\overrightarrow{AB}\| = 6$, $\|\overrightarrow{AC}\| = \sqrt{20} = 2\sqrt{5}$, $\|\overrightarrow{BC}\| = \sqrt{32} = 4\sqrt{2}$.

$\cos\widehat{A} = \dfrac{12}{6 \times 2\sqrt{5}} = \dfrac{1}{\sqrt{5}}$.

$\cos\widehat{B} = \dfrac{\overrightarrow{BA} \cdot \overrightarrow{BC}}{\|\overrightarrow{BA}\| \|\overrightarrow{BC}\|} = \dfrac{24}{6 \times 4\sqrt{2}} = \dfrac{1}{\sqrt{2}}$.

$\overrightarrow{CA}\cdot\overrightarrow{CB}
=\overrightarrow{AC}\cdot\overrightarrow{BC}=8$ : changer le sens des deux vecteurs conserve leur produit scalaire.

$\cos\widehat{C} = \dfrac{\overrightarrow{CA} \cdot \overrightarrow{CB}}{\|\overrightarrow{CA}\| \|\overrightarrow{CB}\|} = \dfrac{8}{2\sqrt{5} \times 4\sqrt{2}} = \dfrac{1}{\sqrt{10}}$.

\item Les trois angles intérieurs appartiennent à $]0,\pi[$. On obtient
\[
\widehat A=\arccos\left(\frac1{\sqrt5}\right)\approx63{,}4^\circ,
\quad \widehat B=\frac\pi4=45^\circ,
\quad \widehat C=\arccos\left(\frac1{\sqrt{10}}\right)\approx71{,}6^\circ.
\]
La somme des valeurs arrondies est $63{,}4^\circ+45^\circ+71{,}6^\circ=180^\circ$.
Ces calculs sont cohérents avec la relation exacte
$\widehat A+\widehat B+\widehat C=\pi$, donnée par le théorème de la somme des angles d'un triangle.
\end{enumerate}
\end{corrige}
